La empresa \textbf{MueblesPro} produce y comercializa tres tipos de productos: \textbf{mesas (M)}, \textbf{sillas (S)} y \textbf{estanterías (E)}, que son transportadas en lotes, y estos pueden estar completos o fraccionados (se puede enviar un lote completo, medio lote, un cuarto de lote, etc.). Por cada lote de estos productos, la empresa obtiene unos ingresos netos de \textbf{20, 15 y 25 unidades monetarias}, respectivamente.

La fábrica dispone de una planta con \textbf{80 horas de producción semanales}. Producir un lote de mesas lleva \textbf{15 horas}, un lote de sillas lleva \textbf{10 horas} y montar un lote de estanterías lleva \textbf{20 horas}.

Además, por contrato con sus distribuidores, la empresa debe suministrar al menos el contenido equivalente a \textbf{6 lotes completos en total}, entre mesas, sillas y estanterías.

El modelo de programación lineal que permite determinar el plan de producción óptimo es el siguiente:

\begin{equation}
\begin{split}
\text{max. } z = 20x_1 + 15x_2 + 25x_3 \\
\text{s.a.:} \\
15x_1 + 10x_2 + 20x_3 \leq 80 \quad \text{(tiempo de montaje disponible)} \\
x_1 + x_2 + x_3 \geq 6 \quad \text{(compromiso de producción total)} \\
x_1, x_2, x_3 \geq 0 \quad \text{(No negatividad)}
\end{split}
\end{equation}

Se pide:

\begin{enumerate}
    \item Obtener el plan de programación de la producción óptimo e indicar los niveles de producción, el beneficio y el cumplimiento de los requisitos descritos. Se puede hacer uso de la información de que, en la solución óptima, se cumple que $x_1=x_3=x_4=h_1=0$.
    \item Describir cuál sería el nuevo plan de producción si fuera necesario entregar al menos 9 lotes en lugar de 6.
    \item Calcular el rango de valores de la contribución unitaria al beneficio de las sillas dentro del cual resulta interesante mantener la base de la solución óptima.
    \item Explicar el significado del criterio del simplex de la variable $x_1$ interpretando explícitamente el significado de las tasas de sustitución de esta variable. La explicación debe referirse al problema de MueblesPro (y no solo a los aspectos matemáticos del modelo de programación lineal planteado).
\end{enumerate}